\documentclass[11pt]{article}

\input{../shared/preamble.tex}

\usepackage{amsmath}
\usepackage{booktabs}
\usepackage{cleveref}

\makeatletter
\renewcommand\paragraph{\@startsection{paragraph}{4}{\z@}%
  {3.25ex \@plus1ex \@minus.2ex}{-1em}%
  {\normalfont\normalsize\itshape}}
\makeatother

\title{Conditioned Cores: A Schur Bridge from Nested Clustered Optimization
to Global Minimum Variance}
\author{Peter Cotton}
\date{September 17, 2026}

\newcommand{\R}{\mathbb{R}}
\newcommand{\ones}{\mathbf{1}}
\DeclareMathOperator{\diag}{diag}
\DeclareMathOperator{\Cov}{Cov}
\DeclareMathOperator{\Var}{Var}

\newtheorem{proposition}{Proposition}
\newtheorem{remark}{Remark}
\newtheorem{definition}{Definition}

\begin{document}
\maketitle

\begin{abstract}
Nested clustered optimization allocates within each cluster from the
cluster's own covariance block and then across the resulting cluster
portfolios. Block inversion says the unconstrained minimum-variance
portfolio has the same two-tier shape, with each block replaced by its Schur
complement against every other asset. We show that under a rank-one model of
cross-cluster dependence, in which each asset reaches other clusters only
through a representative ``core'' of its own cluster, the complement against
everything equals the complement against the other cores, which is a
Vecchia approximation made exact. Damping that
complement by $\gamma\in[0,1]$ gives a bridge with nested clustered
optimization at $\gamma=0$ and the global optimum at $\gamma=1$, with no
inverse larger than a cluster or the number of clusters. The two-stage form
also nests the core-orbital allocation of \textcite{hcoa2026}, from which we
take the word core, and the same model says when that method's strategic use
of the core covariance is accurate.
\end{abstract}

\section{Introduction}

Two-stage cluster methods partition the assets, allocate inside each
cluster, and then allocate across clusters. In nested clustered
optimization, NCO \parencite{lopezdeprado2020}, the inner step optimizes
each cluster on its own covariance block and the outer step optimizes across
the resulting cluster portfolios. The outer problem has one dimension per
cluster, which is the source of the method's numerical stability.

\textcite{hcoa2026} proposes a variant for universes with strong systematic
structure, such as bonds denominated in USD. It is called hierarchical
core-orbital allocation, or HCOA. One representative ``core'' asset is chosen in each
cluster: the medoid, the member most aligned with the cluster's first
principal component, or the member with the largest sum of squared
correlations to the rest. The strategic tier allocates across the cores
with any objective. A tactical rule then redistributes each cluster's
budget among the remaining ``orbital'' members.

In both methods the inner step reads only the cluster's own block. The only
channel by which one cluster's weights respond to another cluster is its
scalar budget.

Schur-complementary allocation \parencite{cotton2024schur} removes the same
blind spot from hierarchical risk parity \parencite{lopezdeprado2016}. At
each split of a bisection tree it replaces a block by its Schur complement
against the sibling block, damped by $\gamma\in[0,1]$, and carries a
companion vector alongside. The result is a bridge with hierarchical risk
parity at $\gamma=0$ and minimum variance at $\gamma=1$.

This note builds the corresponding bridge for a flat partition. The left
end is NCO and the right end is the unconstrained global minimum-variance
portfolio. The obstacle is cost, since the complement of a cluster against
everything else needs an inverse of nearly full size.

The cores remove the obstacle. Under a stated model of cross-cluster
dependence, the other clusters' cores are a sufficient conditioning set,
and no inverse exceeds the size of a cluster or the number of clusters.
Restricting a conditioning set to a few indices is the approximation of
\textcite{vecchia1988}, and the model is the condition under which it loses
nothing. The word core is from \textcite{hcoa2026}, where a core stands in
for its cluster. Here it is a conditioning variable. HCOA is not an endpoint
of the bridge, and \cref{sec:hcoa} says where it sits.

All exactness statements concern the unconstrained, fully invested,
long-short problem. With box or long-only constraints no block
decomposition of the optimum is exact, here or on the tree.

\section{Setup and the two published constructions}

Let $\Sigma\succ0$ be the covariance of $n$ assets partitioned into
clusters $I_1,\dots,I_k$. Cluster $i$ has a core $c_i\in I_i$ and orbitals
$O_i=I_i\setminus\{c_i\}$. Write $C=\{c_1,\dots,c_k\}$ and
$C_{-i}=C\setminus\{c_i\}$. For index sets $X,Y$ write $\Sigma_{XY}$ for the
corresponding block, $\Sigma_{ii}$ for the block of cluster $i$, and $-i$
for the complement of $I_i$.

A two-stage method is specified by two objects per cluster. The
\emph{implemented portfolio} $v_i$, supported on $I_i$ with
$\ones^\top v_i=1$, is where the cluster's budget is spent. The
\emph{strategic proxy} $s_i$ is the portfolio that stands for the cluster
when budgets are set. With $S$ and $V$ the $n\times k$ matrices of proxies
and implemented portfolios, the method is
\begin{equation}
  a \;=\; \operatorname{alloc}\!\big(S^\top\Sigma S\big),
  \qquad w \;=\; Va .
  \label{eq:twostage}
\end{equation}

NCO sets $s_i=v_i$, the minimum-variance portfolio of $\Sigma_{ii}$. The
cluster is priced at the strategic tier by the portfolio it will hold.

HCOA sets $s_i=e_{c_i}$, so that $S^\top\Sigma S=\Sigma_{CC}$, and
$v_i=g_i(\Sigma_{ii})$ for a tactical rule $g_i$ described as
``hierarchical rules or local optimization'' within the
cluster.\footnote{We work from the abstract, the author's announcement,
\url{https://lnkd.in/p/ekyaNa-E}, and indexed excerpts of the full text,
which was not otherwise available to us.} The cluster is priced by its core
and held as something else.

\section{Block inversion}

Fix a vector $u$ and a cluster $i$. Block inversion of $\Sigma$ on the
partition $(I_i,-i)$ gives \parencite[App.~A]{cotton2024schur}
\begin{equation}
  \big(\Sigma^{-1}u\big)_{I_i}
  \;=\; \big(\Sigma^{c}_{ii}\big)^{-1} b_i,
  \label{eq:blockinv}
\end{equation}
where
\begin{equation}
  \Sigma^{c}_{ii} = \Sigma_{ii}-\Sigma_{i,-i}\Sigma_{-i,-i}^{-1}\Sigma_{-i,i},
  \qquad
  b_i = u_{I_i}-\Sigma_{i,-i}\Sigma_{-i,-i}^{-1}u_{-i} .
  \label{eq:complement}
\end{equation}
For a positive-definite $Q$ and a vector $b$ we call $Q^{-1}b$ a
\emph{generalized minimum-variance direction}. Up to scale it minimizes
$x^\top Qx$ subject to $b^\top x=1$, which is ordinary minimum variance only
when $b=\ones$.

With $u=\ones$, the blocks of the global minimum-variance direction
$\Sigma^{-1}\ones$ are therefore generalized minimum-variance directions,
one per cluster, on the cluster's covariance conditional on everything
else. There is no cross-cluster optimizer in the global optimum. The
cluster totals are whatever the stacked directions sum to, and every piece
of cross-cluster information enters through the conditioning.

NCO has the same shape with the conditioning removed. Restoring it costs an
inverse of size $n-|I_i|$ per cluster, which is what a two-stage method
exists to avoid.

\section{The orbital model}

\begin{definition}[Orbital model]
For every cluster $j$, the residual of the orbitals $O_j$ after linear
regression on their own core is uncorrelated with every asset outside
$I_j$. Equivalently the cross-cluster residual
\begin{equation}
  R_j(c_j) \;=\; \Sigma_{O_j,-j}-\Sigma_{O_j c_j}\,\sigma_{c_j}^{-2}\,\Sigma_{c_j,-j}
  \label{eq:residual}
\end{equation}
is zero.
\end{definition}

An orbital may be correlated with anything, but its correlation with other
clusters passes entirely through its core. Each cross-cluster block of
$\Sigma$ then has rank one.

This is a strong assumption, and clustering alone does not imply it. It
asks that orbital residuals be uncorrelated both with the other cores and
with the other clusters' orbital residuals.

\begin{proposition}[Cores are sufficient]
\label{prop:sufficient}
Under the orbital model, for each cluster $i$ and every $u$,
\begin{equation}
  \Sigma^{c}_{ii}
  \;=\; \Sigma_{ii}-\Sigma_{i,C_{-i}}\,\Sigma_{C_{-i}C_{-i}}^{-1}\,\Sigma_{C_{-i},i},
  \qquad
  b_i \;=\; u_{I_i}-\Sigma_{i,C_{-i}}\,\Sigma_{C_{-i}C_{-i}}^{-1}\,u_{C_{-i}} .
  \label{eq:cheap}
\end{equation}
\end{proposition}

\begin{proof}
The matrix $\Sigma_{i,-i}\Sigma_{-i,-i}^{-1}$ holds the coefficients of the
linear regression of cluster $i$ on all other assets. Split the regressors
into the other cores $C_{-i}$ and the other orbitals $O_{-i}$, and write
$O_{-i}=\beta\,C_{-i}+\varepsilon$ for the block-diagonal regression of
each orbital set on its own core. The map from $(C_{-i},O_{-i})$ to
$(C_{-i},\varepsilon)$ is invertible, so the regression may be run on the
latter pair. Under the orbital model $\varepsilon$ is uncorrelated with
cluster $i$ and with $C_{-i}$, so its coefficients vanish and the
coefficients on $C_{-i}$ are $\Sigma_{i,C_{-i}}\Sigma_{C_{-i}C_{-i}}^{-1}$.
Substituting into \eqref{eq:complement} gives \eqref{eq:cheap}, with
$u_{-i}$ contributing only through its $C_{-i}$ entries.
\end{proof}

Conditioning cluster $i$ on the other cores is the same as conditioning it
on everything. The inverse in \eqref{eq:cheap} is $(k-1)\times(k-1)$.

\begin{remark}[A Vecchia approximation]
Replacing the conditioning set $-i$ in \eqref{eq:complement} by the index
set $C_{-i}$ is a Vecchia approximation \parencite{vecchia1988}. The
conditionals there are the same Schur complements, and the approximation is
a choice of which indices to condition on. In the framework of
\textcite{katzfuss2021} most Gaussian-process approximations are choices of
this kind. The orbital model is the conditional independence under which
this particular choice is exact.
\end{remark}

\begin{remark}[Testing the model and choosing cores]
Regressing each orbital on its own core and the other cores, and checking
that the other-core coefficients vanish, does not test the model. Orbital
residuals can be orthogonal to every core and still be correlated with each
other across clusters, and the sufficiency of the cores then fails. The
direct diagnostic is the size of $R_j(c_j)$ in \eqref{eq:residual},
computed from the sample covariance before any portfolio is built. The same
quantity gives a core-selection rule aligned with
\cref{prop:sufficient}: choose $c_j$ to minimize $\|R_j(c)\|$ over
$c\in I_j$. The medoid, principal-component and squared-correlation
criteria of \textcite{hcoa2026} read only within-cluster structure.
\end{remark}

\begin{remark}[A symmetric factor model]
Suppose instead that every member of cluster $j$ loads on a latent cluster
factor, $r_m=\lambda_m f_j+\eta_m$, with idiosyncratic terms uncorrelated
with each other and with the factors. Then
$R_j(c)=(1-\rho_c)\,\lambda_{O_j}\Cov(f_j,r_{-j})$, where
$\rho_c=\lambda_c^2\Var(f_j)/\sigma_c^2$ is the share of the core's variance
explained by the factor. The orbital model holds to the extent the core is
the factor. In this model the within-cluster criteria that favour a high
$\rho_c$ are approximately the right ones.
\end{remark}

\section{The bridge}

Following \textcite{cotton2024schur}, damp the conditioning. For
$\gamma\in[0,1]$ define the \emph{conditioned pair} of cluster $i$,
\begin{equation}
\begin{aligned}
  Q_i(\gamma) &= \Sigma_{ii}-\gamma\,\Sigma_{i,C_{-i}}\Sigma_{C_{-i}C_{-i}}^{-1}\Sigma_{C_{-i},i},\\
  b_i(\gamma) &= u_{I_i}-\gamma\,\Sigma_{i,C_{-i}}\Sigma_{C_{-i}C_{-i}}^{-1}u_{C_{-i}},
\end{aligned}
  \label{eq:pair}
\end{equation}
and the direction $d_i(\gamma)=Q_i(\gamma)^{-1}b_i(\gamma)$. When
$\ones^\top d_i(\gamma)\neq0$ set $v_i=d_i/\ones^\top d_i$, take $s_i=v_i$
in \eqref{eq:twostage}, and let the strategic tier solve the same problem
on the span of $V$, which gives
\begin{equation}
  a \;\propto\; \big(V^\top\Sigma V\big)^{-1}V^\top u,
  \qquad w = Va .
  \label{eq:top}
\end{equation}
The induced strategic inputs are the $k\times k$ covariance $V^\top\Sigma V$
and the $k$-vector $V^\top u$, which is $\ones$ when $u=\ones$.

\begin{proposition}[The two ends of the bridge]
\label{prop:bridge}
\leavevmode
\begin{enumerate}
\item At $\gamma=0$ with $u=\ones$ the recipe is NCO with minimum variance
  at both tiers.
\item Under the orbital model the stacked directions $d_i(1)$ are the blocks
  of $\Sigma^{-1}u$.
\item If in addition every $\ones^\top d_i(1)$ is nonzero, then $w$ in
  \eqref{eq:top} is proportional to $\Sigma^{-1}u$.
\end{enumerate}
\end{proposition}

\begin{proof}
At $\gamma=0$ the pair is $(\Sigma_{ii},\ones)$, so $v_i$ is the
minimum-variance portfolio of the block and \eqref{eq:top} is minimum
variance across the cluster portfolios. The second claim is
\eqref{eq:blockinv} with \cref{prop:sufficient}. For the third,
$w^\star=\Sigma^{-1}u$ lies in the column span of $V$, say
$w^\star=V\alpha$. Then
$(V^\top\Sigma V)^{-1}V^\top u
=(V^\top\Sigma V)^{-1}V^\top\Sigma V\alpha=\alpha$.
\end{proof}

\begin{table}[ht]
\centering
\begin{tabular}{@{}ll@{}}
\toprule
$\gamma$ & Method \\
\midrule
$0$ & NCO \parencite{lopezdeprado2020} \\
$(0,1)$ & conditioned cores \\
$1$ & unconstrained global minimum variance, under the orbital model \\
\bottomrule
\end{tabular}
\caption{The bridge with $u=\ones$ and minimum variance at both tiers.}
\label{tab:bridge}
\end{table}

\paragraph{Scope.} The right end is the unconstrained, fully invested,
long-short solution. The third claim also needs every cluster total to be
nonzero. A global solution can hold a self-financing position inside a
cluster with zero net weight, and then no cluster budget times a normalized
cluster portfolio represents it. The second claim is unaffected, because it
stacks unnormalized directions.

\begin{remark}[Other objectives]
\Cref{prop:bridge} is stated for a direction $\Sigma^{-1}u$. With $u=\mu$
this is the maximum Sharpe direction, normalized afterwards when
$\ones^\top\Sigma^{-1}\mu\neq0$. Mean--variance problems with a budget
constraint have solutions that combine $\Sigma^{-1}\mu$ and
$\Sigma^{-1}\ones$, and each part is covered separately. Risk parity is not
of this form. The recipe still runs with a risk-parity strategic tier, but
nothing is exact at $\gamma=1$.
\end{remark}

\begin{remark}[Rules that read only a covariance]
A tactical rule that cannot take a constraint vector can be applied after
a change of variables. With $y=b_i\circ x$ the constraint $b_i^\top x=1$
becomes $\ones^\top y=1$ and the covariance becomes
$M=\diag(b_i)^{-1}Q_i\diag(b_i)^{-1}$. Apply the rule to $M$ and map back by
$x=y/b_i$, coordinate by coordinate. For minimum variance this returns
$Q_i^{-1}b_i$. It requires $b_i$ to have no zero entry. Applying minimum
variance to $M$ without mapping back returns $\diag(b_i)Q_i^{-1}b_i$, which
is a different portfolio.
\end{remark}

\begin{remark}[Relation to the tree]
Schur complements compose: the complement of a complement is the complement
against the union. So at $\gamma=1$ the recursion of
\textcite{cotton2024schur}, run down any tree whose leaves are the clusters
and stopped there, hands each cluster the pair \eqref{eq:complement}. The
flat construction and the iterated one agree at the right end. They differ
in the interior, because the tree damps at every split and the damping
compounds along the path, while \eqref{eq:pair} damps each cluster once
against all the other cores.
\end{remark}

\begin{remark}[What $\gamma$ means]
The cross-core regression coefficients in \eqref{eq:pair} are estimated.
\textcite{cotton2026bridge} show that on the bisection tree, with
estimation error, the out-of-sample-optimal $\gamma$ generically lies
strictly inside $(0,1)$. We conjecture the same here, since the cross-core
block is the only object $\gamma$ touches, but we do not prove it. Core
covariances in a bond universe are highly correlated, so the inverse in
\eqref{eq:pair} is where we would expect damping to matter.
\end{remark}

\section{Where HCOA sits}
\label{sec:hcoa}

The two-stage form \eqref{eq:twostage} nests all three methods.
\Cref{tab:nest} lists the proxy and the implemented portfolio for each.

\begin{table}[ht]
\centering
\begin{tabular}{@{}lll@{}}
\toprule
Method & Strategic proxy $s_i$ & Implemented $v_i$ \\
\midrule
NCO & minimum variance on $\Sigma_{ii}$ & same as $s_i$ \\
HCOA & the core, $e_{c_i}$ & $g_i(\Sigma_{ii})$ \\
Conditioned cores & $d_i(\gamma)$, normalized & same as $s_i$ \\
\bottomrule
\end{tabular}
\caption{Three instances of the two-stage form \eqref{eq:twostage}.}
\label{tab:nest}
\end{table}

HCOA differs from NCO and from the bridge on a second axis. Its strategic
proxy is the core while its implemented portfolio is something else.
Changing $\gamma$ alone does not reach it, and setting $v_i=e_{c_i}$ in the
bridge would hold only the cores.

The orbital model says what the core covariance is a proxy for.

\begin{proposition}[The core covariance is the systematic part]
\label{prop:proxy}
Under the orbital model, let $v_i$ be any portfolios supported on their
clusters. Write $\beta_i=\Sigma_{O_ic_i}/\sigma_{c_i}^2$ for the orbital
betas, $E_i=\Sigma_{O_iO_i}-\beta_i\beta_i^\top\sigma_{c_i}^2$ for the
residual covariance, and $\kappa_i=v_{i,c_i}+\beta_i^\top v_{i,O_i}$ for
the core beta of the cluster portfolio. Then
\begin{equation}
  V^\top\Sigma V \;=\; K\,\Sigma_{CC}\,K + D,
  \qquad K=\diag(\kappa_i),
  \quad D=\diag\!\big(v_{i,O_i}^\top E_i\,v_{i,O_i}\big).
  \label{eq:proxy}
\end{equation}
\end{proposition}

\begin{proof}
Under the orbital model $\Sigma_{O_i,l}=\beta_i\Sigma_{c_i,l}$ for every
asset $l$ outside $I_i$. Applying this on both sides gives
$v_i^\top\Sigma_{I_iI_l}v_l=\kappa_i\kappa_l\Sigma_{c_ic_l}$ for $i\neq l$.
On the diagonal, the core is uncorrelated with the orbital residual, so
$v_i^\top\Sigma_{ii}v_i=\kappa_i^2\sigma_{c_i}^2+v_{i,O_i}^\top E_iv_{i,O_i}$.
\end{proof}

The consistent strategic matrix is the core covariance scaled by the
cluster portfolios' core betas, plus their residual variances on the
diagonal. HCOA's strategic tier is \eqref{eq:proxy} with $K=I$ and $D=0$.

It is accurate when the implemented cluster portfolios have unit beta to
their cores and little residual variance, which describes tight,
homogeneous clusters. The correction is still $k\times k$ and needs nothing
beyond $\beta_i$ and $E_i$, which the tactical tier already has.

\section{Cost}

No step forms an $n\times n$ inverse. The largest linear solve has dimension
$\max(k,\max_i|I_i|)$, as in NCO: the strategic tier is $k\times k$ and each
conditioning inverse is $(k-1)\times(k-1)$. The bridge adds work that NCO does not
do: the cross-block products $\Sigma_{i,C_{-i}}$, and a leave-one-core-out
inverse per cluster, obtainable from $\Sigma_{CC}^{-1}$ by a rank-one
downdate.

If the orbital model fails, \eqref{eq:cheap} approximates the full
complement in \eqref{eq:complement}, and the recipe is a Schur allocation
with a truncated conditioning set. The size of $R_j$ in
\eqref{eq:residual} measures the truncation.

\section{Closing}

The contribution is deliberately small.\footnote{\texttt{verify\_conditioned\_cores.py}
next to the manuscript checks \cref{prop:sufficient,prop:bridge,prop:proxy},
the two remarks on diagnostics and on covariance-only rules, and the
symmetric-factor formula, on random covariances to $10^{-10}$. It also
confirms that the exactness claims fail when the orbital model is violated.
The repository's test suite runs it.} The identity is block inversion, the
restricted conditioning set is Vecchia's, the damping is from
\textcite{cotton2024schur}, and the word core is from \textcite{hcoa2026}. What is added is the observation that under a
rank-one cross-cluster model the cores are a sufficient conditioning set,
which makes the bridge from NCO to the global optimum cheap.

\Cref{prop:sufficient} may also serve Schur-complementary allocation on the
tree, where the complement against a sibling block could be taken against
that block's cores. We leave that, and an empirical comparison on a bond
universe, to future work.

\printbibliography

\end{document}
